\documentclass[10pt]{article}
\usepackage[a4paper,margin=0.6in]{geometry}
\usepackage{amsmath,amssymb,amsfonts,mathtools}
\setlength{\jot}{4pt}
\usepackage{xcolor}
\setlength{\parindent}{0pt}
\setlength{\parskip}{5pt}
\newcommand{\Leps}{L_{\epsilon}}
\newcommand{\KK}{\mathcal K(x-y)\!\cdot\!\mathcal K(x'-y')}
\newcommand{\Wbr}{\left[U^{ce'}(y') - U^{ce'}(x')\right] \left[U^{ce}(y) - U^{ce}(x)\right] \rho^{e'}(x')\, \rho^{e}(x)}
\title{Group II, Row IV: extraction of the BFKL, JIMWLK and finite\\ (DGLAP $+$ RC) terms}
\author{}
\date{}
\begin{document}
\maketitle
\vspace{-2.5em}

\section{Check of the starting expression}
The starting expression is
\begin{align}\label{eq:start}
     \frac{d^{3}N_{\rm NLO}^{\rm}}{d^{3}k}  & =  \frac{1}{(2\pi)^{3}}\,\frac{g^{4}\,d_{\perp}}{4\pi^{3}}\,\frac{N_{c}}{k^{+}}\,\int_{\frac{k^{+}}{k^{+}+\Lambda}}^{\frac{k^{+}}{\vee}}\, \frac{d\xi}{2\pi}\int_{x,y,y^{\prime},x'}\,e^{-ik\,\left(y'-y\right)\,\frac{1}{\xi}}\,\frac{(x-y)\cdot (x'-y')}{(x-y)^{2}\,(x'-y')^{2}}\,\bigg[\xi\,(1-\xi)+\frac{(1-\xi)}{\xi}+\frac{\xi}{(1-\xi)}\bigg]\,\frac{1}{\xi^{2}}\notag \\ & \left[\frac{1}{\epsilon} +\gamma_{E} +\frac{1}{2}\,\log{\left[\frac{(y'-y)^{2}\,\mu^{2}_{\overline{MS}}}{4}\right]}\right]\,\Wbr\notag \\ & +\,\frac{1}{(2\pi)^{3}}\,\frac{g^{4}}{4\pi^{3}}\,\frac{N_{c}}{k^{+}}\,\int_{\frac{k^{+}}{k^{+}+\Lambda}}^{\frac{k^{+}}{\vee}}\, \frac{d\xi}{2\pi}\,\int_{x,y,y^{\prime},x'}\,e^{-ik\,\left(y'-y\right)\,\frac{1}{\xi}}\,\frac{1}{\xi^{2}}\Bigg[\frac{(x-y)\cdot (y'-y)\, (x'-y')\cdot (y'-y)}{(x-y)^{2}\,(x'-y')^{2}\,(y'-y)^{2}}\,d_{\perp}\,\xi(1-\xi)\notag \\ & + \frac{(x-y)\cdot (x'-y')}{(x-y)^{2}\,(x'-y')^{2}}\,\left[\frac{(1-\xi)}{\xi}+\frac{\xi}{(1-\xi)}\right]\Bigg]\,\times \Wbr\,.
\end{align}
\textbf{It is correct}: it coincides exactly with the reduction of the Row-IV master expression
(colour algebra $U^{d'b}(z)U^{db}(z)=\delta^{d'd}$, $f^{c'da}f^{cda}=N_c\delta^{c'c}$;
$\delta$-function integration over $w,w'$; contraction of the two tensor brackets over $i,j$,
in which the antisymmetric $\frac{2}{k^{+}}(\delta_{mk'}\delta_{m'k}-\delta_{m'k'}\delta_{km})$
channel cancels; the dimensionally regularized $z$-integral
$\int_{z}\mathcal K^{m}(y-z)\mathcal K^{m'}(y'-z)=-\frac{\pi}{2}\Leps\delta^{mm'}-\pi\frac{r^{m}r^{m'}}{r^{2}}$,
$r=y'-y$, $\Leps\equiv\frac{2}{\epsilon}+2\gamma_{E}+\log\frac{r^{2}\mu^{2}_{\overline{MS}}}{4}$;
and the substitution $\xi=\frac{k^{+}}{p^{+}+k^{+}}$), as derived step by step in the companion note
\emph{``Group II, Row IV: reduction of the master expression''}. Two features were checked
explicitly: (i) the $\xi$-limits are written \emph{descending}, from
$\xi_{\max}=\frac{k^{+}}{k^{+}+\Lambda}$ down to $\bar\xi=\frac{k^{+}}{\vee}$, which absorbs the sign
of $dp^{+}=-\frac{k^{+}d\xi}{\xi^{2}}$, so all the printed signs are right; (ii) the placement of
$d_{\perp}$ is right: it multiplies the full $\mathcal P$-combination in the $\Leps$ term, but only
the azimuthal kernel in the second term --- the dipole piece
$\big[\frac{1-\xi}{\xi}+\frac{\xi}{1-\xi}\big]$ carries no $d_{\perp}$, since it originates from the
$-\pi\,\delta_{mm'}\frac{r^{m}r^{m'}}{r^{2}}=-\pi$ trace of the $z$-integral.

\section{Notation and first regrouping}
Write $C\equiv\frac{1}{(2\pi)^{3}}\frac{g^{4}N_{c}}{8\pi^{4}}\frac{1}{k^{+}}
=\frac{g^{2}}{4\pi^{2}}\frac{N_{c}}{(2\pi)^{3}}\frac{g^{2}}{2\pi^{2}}\frac{1}{k^{+}}$,
$f(\xi)=e^{-ik(y'-y)/\xi}$, $\hat p(\xi)=\xi(1-\xi)+\frac{1-\xi}{\xi}+\frac{\xi}{1-\xi}$,
and keep the colour bracket of \eqref{eq:start} implicit. Flipping the limits to ascending
($\int_{\xi_{\max}}^{\bar\xi}=-\int_{\bar\xi}^{\xi_{\max}}$), Eq.~\eqref{eq:start} reads
\begin{align}\label{eq:ascending}
\frac{d^{3}N_{\rm NLO}}{d^{3}k}
=-\,C\!\int_{\bar\xi}^{\xi_{\max}}\!\!d\xi\,\Big\{\frac{d_{\perp}}{2}\,\frac{\hat p(\xi)}{\xi^{2}}\,\Leps
+\frac{1}{\xi^{2}}\Big[\frac{\xi}{1-\xi}+\frac{1-\xi}{\xi}\Big]\Big\}\,\KK\,f(\xi)
-C\,d_{\perp}\!\int_{\bar\xi}^{\xi_{\max}}\!\!d\xi\,\frac{1-\xi}{\xi}\,
\frac{[\mathcal K(x-y)\!\cdot\! r][\mathcal K(x'-y')\!\cdot\! r]}{r^{2}}\,f(\xi)\,.
\end{align}
Using the exact identity
\begin{align}\label{eq:phatid}
\frac{\hat p(\xi)}{\xi^{2}}-\frac{1-\xi}{\xi}
=\frac{1}{\xi^{2}}\Big[\frac{\xi}{1-\xi}+\frac{1-\xi}{\xi}\Big]\,,
\end{align}
the $\Leps$-free dipole piece is rewritten as $\hat p/\xi^{2}$ minus $\frac{1-\xi}{\xi}$, giving
\begin{align}\label{eq:regrouped}
\frac{d^{3}N_{\rm NLO}}{d^{3}k}
=-\,C\!\int_{\bar\xi}^{\xi_{\max}}\!\!d\xi\,\frac{\hat p(\xi)}{\xi^{2}}
\underbrace{\Big[\frac{d_{\perp}}{2}\Leps+1\Big]}_{=\;\Leps\;+\;O(\epsilon)}\KK f(\xi)
\;+\;C\!\int_{\bar\xi}^{1}\!\!d\xi\,\frac{1-\xi}{\xi}
\Big[\KK-d_{\perp}\,\frac{(\mathcal K\!\cdot\! r)(\mathcal K'\!\cdot\! r)}{r^{2}}\Big]f(\xi)\,,
\end{align}
where $\frac{d_{\perp}}{2}\Leps+1=\frac{d_{\perp}}{\epsilon}+d_{\perp}\gamma_{E}
+\frac{d_{\perp}}{2}\log\frac{r^{2}\mu^{2}}{4}+1=\Leps$ for $d_{\perp}=2-\epsilon$
(the $+1$ from $-\frac{\epsilon}{2}\cdot\frac{2}{\epsilon}$ cancels the explicit $+1$), and the
upper limit of the regular second integral has been extended to $1$ at a cost $O(\Lambda)$.

\section{Endpoint extraction: the JIMWLK term}
A $\log\Lambda$ can only come from the upper limit
$\xi_{\max}=\frac{k^{+}}{k^{+}+\Lambda}\simeq1-\frac{\Lambda}{k^{+}}$. Of the three pieces of
$\hat p(\xi)$, $\xi(1-\xi)$ and $\frac{1-\xi}{\xi}$ are regular at $\xi=1$, so for them
$\xi_{\max}\to1$ at a cost $O(\Lambda)$. The entire $\Lambda$-dependence of \eqref{eq:regrouped}
therefore sits in the single channel $\frac{\xi}{1-\xi}$, which is singular at $\xi\to1$ --- and
$\xi\to1$ means $p^{+}\to0$: \emph{the unmeasured gluon becomes soft}. The logarithm regulated by
$\Lambda$ at this endpoint is the target-side rapidity logarithm, i.e.\ the JIMWLK piece.
Isolating that channel and writing its integrand as $h(\xi)/(1-\xi)$ with
\begin{align}
h(\xi)\equiv\Big[\frac{d_{\perp}}{2}\Leps+1\Big]\,\frac{f(\xi)}{\xi}\,,
\end{align}
we add and subtract the endpoint value $h(1)$ in the numerator:
\begin{align}\label{eq:addsub}
\int_{\bar\xi}^{\xi_{\max}}\!\!d\xi\,\frac{h(\xi)}{1-\xi}
=\int_{\bar\xi}^{\xi_{\max}}\!\!d\xi\,\frac{h(\xi)-h(1)}{1-\xi}
\;+\;h(1)\int_{\bar\xi}^{\xi_{\max}}\frac{d\xi}{1-\xi}\,.
\end{align}
In the first piece the numerator vanishes linearly at $\xi=1$, so the integrand is finite there
and $\xi_{\max}\to1$ at a cost $O(\Lambda)$: this finite remainder is what the locally subtracted
$\frac{\xi}{(1-\xi)_{+}}$ denotes. The second piece is elementary:
\begin{align}\label{eq:elementary}
\int_{\bar\xi}^{\xi_{\max}}\frac{d\xi}{1-\xi}
=\log\frac{1-\bar\xi}{1-\xi_{\max}}
=\log\frac{\big(1-\frac{k^{+}}{\vee}\big)(k^{+}+\Lambda)}{\Lambda}
=\log\frac{(\vee-k^{+})\,k^{+}}{\vee\,\Lambda}+O(\Lambda)\,,
\end{align}
\textbf{Sign check of \eqref{eq:elementary}.} The antiderivative of $\frac{1}{1-\xi}$ is
$-\log(1-\xi)$, so substituting the limits in the standard order, \emph{upper minus lower},
\begin{align*}
\int_{\bar\xi}^{\xi_{\max}}\frac{d\xi}{1-\xi}
=\Big[-\log(1-\xi)\Big]_{\xi=\bar\xi}^{\,\xi=\xi_{\max}}
=-\log(1-\xi_{\max})-\big[-\log(1-\bar\xi)\big]
=\log\frac{1-\bar\xi}{1-\xi_{\max}}\,:
\end{align*}
it is the minus sign of the antiderivative that puts the \emph{lower} limit in the numerator of
the ratio; the limits themselves are substituted the usual way around. Consistency check: the
integrand $\frac{1}{1-\xi}$ is positive on the whole range and the limits are ascending
($\bar\xi<\xi_{\max}$), so the integral must be positive --- and indeed
$1-\bar\xi=\frac{\vee-k^{+}}{\vee}=O(1)$ while $1-\xi_{\max}=\frac{\Lambda}{k^{+}+\Lambda}\to0$,
so \eqref{eq:elementary} is large and positive. (The descending limits of the original expression
\eqref{eq:start} were already traded for the overall minus sign of \eqref{eq:ascending}; no second
sign flip occurs here.)

\textbf{The endpoint value $h(1)$.} The only $\xi$-dependence left in $h$ is the explicit
$\frac1\xi$ and the phase $f(\xi)=e^{-ik(y'-y)/\xi}$: the transverse kernels
$\mathcal K(x-y)\cdot\mathcal K(x'-y')$ and the bracket $\frac{d_{\perp}}{2}\Leps+1$ carry no
$\xi$ at all. The $\frac1\xi$ in the phase is purely kinematic: in the master expression the
phase is $e^{-ik(y'-y)\frac{p^{+}+k^{+}}{k^{+}}}$, and
$\frac{p^{+}+k^{+}}{k^{+}}=\frac1\xi$. At the endpoint $\xi=1$, i.e.\ $p^{+}=0$, the unmeasured
gluon carries zero longitudinal momentum, so the measured gluon retains the full leading-order
kinematics: the phase collapses to the LO phase, $f(1)=e^{-ik(y'-y)}$, and the explicit
$\frac1\xi\to1$. Hence
\begin{align}\label{eq:endpoint}
h(1)=\Big[\frac{d_{\perp}}{2}\Leps+1\Big]\,e^{-ik(y'-y)}\,,
\qquad
\log\frac{(\vee-k^{+})k^{+}}{\vee\Lambda}
={\color{blue}\log\frac{k^{+}}{\Lambda}}+\log\Big(1-\frac{k^{+}}{\vee}\Big)\,.
\end{align}
The coefficient of the endpoint logarithm is thus the $\Leps$-dressed LO dipole --- the soft
unmeasured gluon is the JIMWLK gluon. This channel enters \eqref{eq:regrouped} with the prefactor
$-C$, so its endpoint piece is $-C\,h(1)\times\eqref{eq:elementary}$; pulling the minus sign into
the bracket, $-h(1)=\big[-\frac{d_{\perp}}{2}\Leps-1\big]e^{-ik(y'-y)}$, which is where the
overall $+$ sign and the all-negative bracket of the box below both come from (one minus sign,
used once --- it originates from flipping the descending limits of \eqref{eq:start}, not from the
limit substitution in \eqref{eq:elementary}). The $\log\frac{k^{+}}{\Lambda}$ part is the
\textbf{JIMWLK contribution of this row},
\begin{align}\label{eq:jimwlk}
\boxed{\;\frac{d^{3}N^{\rm JIMWLK}_{\rm G\text{-}II,\,R\,IV}}{d^{3}k}
= + \frac{1}{(2\pi)^{3}}\,\frac{g^{4}N_{c}}{8\pi^{4}}\, \frac{1}{k^{+}}\,{\color{blue}\log\Big(\frac{k^{+}}{\Lambda}\Big)}\int_{x,y,x',y'}e^{-ik\cdot (y'-y)}\,\KK\,
\Big[-\frac{d_{\perp}}{\epsilon}-d_{\perp}\gamma_{E}-\frac{d_{\perp}}{2}\log\frac{(y'-y)^{2}\mu^{2}_{\overline{MS}}}{4}-1\Big]\Wbr\;}
\end{align}
(the bracket is $-\Leps$), which is moved to the two-$\rho$ evolution (JIMWLK) file, where it
appears verbatim. The $\log(1-\frac{k^{+}}{\vee})$ part of \eqref{eq:endpoint} stays here.

\section{The standard plus prescription: DGLAP and RC terms}
With the \emph{standard} plus distribution on $[0,1]$,
\begin{align}\label{eq:stdplus}
\int_{a}^{1}d\xi\,\frac{h(\xi)}{(1-\xi)_{+}}=\int_{a}^{1}d\xi\,\frac{h(\xi)-h(1)}{1-\xi}+h(1)\,\log(1-a)\,,\qquad a=\bar\xi=\frac{k^{+}}{\vee}\,,
\end{align}
the locally subtracted integral of \eqref{eq:endpoint} and the leftover
$h(1)\log(1-\frac{k^{+}}{\vee})$ term recombine into the single integral
$\int_{\bar\xi}^{1}d\xi\,h(\xi)/(1-\xi)_{+}$: the two groupings differ only by where the finite
piece $h(1)\log(1-\bar\xi)$ is placed. This conversion changes nothing, but it is required for
\begin{align}\label{eq:pgg}
\mathcal{P}_{gg,+}(\xi) = 2N_{c}\,\Big[\xi(1-\xi)+\frac{1-\xi}{\xi}+\frac{\xi}{(1-\xi)_{+}}\Big]\,,
\qquad
\mathcal{P}_{gg}(\xi) = \mathcal{P}_{gg,+}(\xi) + \frac{\beta_{0}}{2}\,\delta(1-\xi)\,,\qquad
\beta_{0}=\frac{11N_{c}}{3}\ (\text{pure glue})
\end{align}
to be the \emph{standard} splitting function. Using $C=\frac{g^{2}}{8\pi^{2}}\cdot2N_{c}\cdot
\frac{1}{(2\pi)^{3}}\frac{g^{2}}{2\pi^{2}}\frac{1}{k^{+}}$ and defining
\begin{align}
\frac{d^{3}N_{\xi\,\rm LO}}{d^{3}k} = \frac{1}{(2\pi)^{3}}\, \frac{g^{2}}{2\pi^{2}}\, \frac{1}{k^{+}}\,\int_{x,y,x',y'}\, e^{-ik\cdot (y'-y)\,\frac{1}{\xi}}\, \KK\;\Wbr\,,
\end{align}
the first term of \eqref{eq:regrouped} minus the JIMWLK term becomes
$-\frac{g^{2}}{8\pi^{2}}\int_{\bar\xi}^{1}\frac{d\xi}{\xi^{2}}\,\mathcal P_{gg,+}(\xi)\,\Leps\,
\frac{d^{3}N_{\xi\rm LO}}{d^{3}k}$. Finally, adding and subtracting the
$\frac{\beta_{0}}{2}\delta(1-\xi)$ piece,
\begin{align}\label{eq:dglaprc}
-\frac{g^{2}}{8\pi^{2}}\int_{\bar\xi}^{1}\frac{d\xi}{\xi^{2}}\,\mathcal P_{gg,+}\,\Leps\,\frac{d^{3}N_{\xi\rm LO}}{d^{3}k}
=\underbrace{-\frac{g^{2}}{8\pi^{2}}\int_{\bar\xi}^{1}\frac{d\xi}{\xi^{2}}\,\mathcal P_{gg}(\xi)\,\Leps\,\frac{d^{3}N_{\xi\rm LO}}{d^{3}k}}_{\textstyle\rm DGLAP}
\;+\;\underbrace{\frac{g^{2}\beta_{0}}{16\pi^{2}}\,\Leps\,\frac{d^{3}N_{\rm LO}}{d^{3}k}}_{\textstyle\rm RC}\,,
\end{align}
which makes the running-coupling (RC) term explicit; the two added pieces cancel identically, so
this is an exact add-and-subtract.

\section{The soft subtraction: the BFKL term}
\textbf{Step 1: where the second rapidity logarithm hides.}
The DGLAP term \eqref{eq:dglaprc} contains the soft channel $2N_{c}\frac{1-\xi}{\xi}$ of
$\mathcal P_{gg}$, whose lower endpoint is $\bar\xi=\frac{k^{+}}{\vee}$. Now $\xi\to0$ means
$p^{+}\to\vee$: it is the \emph{unmeasured} gluon that carries almost all of the projectile
longitudinal momentum, and the \emph{measured} gluon is the soft one --- the mirror image of the
JIMWLK endpoint of Sec.~3. This is the projectile-side rapidity region, so the associated
logarithm is $\log\frac{\vee}{k^{+}}$ and must be BFKL. In position space the $\xi\to0$ region is
the collapse $y=\xi w+(1-\xi)z\to z$, $y'=\xi w'+(1-\xi)z\to z$ of the parent positions onto the
sibling at fixed $w,w'$. The danger is that in \eqref{eq:dglaprc} this channel multiplies
$\Leps$: left as it stands, it would generate
$\log\frac{\vee}{k^{+}}\times\big(\frac{2}{\epsilon}+\log\mu^{2}+\dots\big)$, and a rapidity
logarithm with an ultraviolet coefficient is not physical --- the BFKL kernel is finite in $d=2$
and $\mu$-independent. The cure is to subtract from the soft channel an expression with the
\emph{same} $\xi\to0$ limit that can be computed exactly, and to add its value back.

\textbf{Step 2: the subtraction term.} The $\xi\to0$ limit is subtracted \emph{exactly} with
\begin{align}\label{eq:sub}
\frac{d^{3}N^{\rm sub}_{\xi}}{d^{3}k} = \frac{1}{(2\pi)^{3}}\, \frac{g^{2}}{2\pi^{2}}\, \frac{1}{k^{+}}\,\int_{x,y,x',y',z}\, e^{-ik\cdot (y'-y)\,\frac{1}{\xi}}\;\frac{1}{\pi}\,K(y,y',z)\,K(x,x',z)
\left[U^{ce'}(z) - U^{ce'}(x')\right]\left[U^{ce}(z) - U^{ce}(x)\right]\rho^{e'}(x')\,\rho^{e}(x)\,,
\end{align}
$K(a,b,z)=\frac{(a-z)\cdot(b-z)}{(a-z)^{2}(b-z)^{2}}$. This is the real-emission structure with
the parent sitting at the common point $z$: the kernels $K(y,y',z)\,K(x,x',z)$ are the two
soft-emission (dipole) kernels of amplitude and conjugate, and the Wilson-line factor is taken at
coincident points, $W(z,z)$ --- exactly the building blocks of the real BFKL kernel. It is IR and
UV finite in $d=2$, contains no $\mu$, and is symmetric under $(x,y)\leftrightarrow(x',y')$.

\textbf{Step 3: the combination is free of $\frac{2}{\epsilon}$ and $\mu$.} Writing
$W(u,u')=[U^{ce'}(u')-U^{ce'}(x')][U^{ce}(u)-U^{ce}(x)]\rho^{e'}(x')\rho^{e}(x)$ and using the
dimensionally regularized identity $\int_{z}^{(d=2-\epsilon)}K(y,y',z)=-\pi\Leps$ to put the
$\Leps$ of the soft channel itself under a $z$-integral, the per-$\xi$ combination
\begin{align}\label{eq:combo}
2N_c\frac{1-\xi}{\xi}\Big[\Leps\,\frac{d^{3}N_{\xi\,\rm LO}}{d^{3}k}+\frac{d^{3}N^{\rm sub}_{\xi}}{d^{3}k}\Big]
=-\frac{2N_c(1-\xi)}{\xi}\,\frac{1}{(2\pi)^{3}}\frac{g^{2}}{2\pi^{2}}\frac{1}{k^{+}}
\int e^{-ik(y'-y)\frac1\xi}\,\frac{K(y,y',z)}{\pi}\Big[\KK\,W(y,y')-K(x,x',z)\,W(z,z)\Big]
\end{align}
puts both pieces under one $z$-integral. To see what this buys, change variables at fixed $z$,
\begin{align}\label{eq:rescale}
y=z+\xi\,(w-z)\,,\qquad y'=z+\xi\,(w'-z)\,,
\end{align}
which is an exact substitution, not an approximation. Under it
\begin{align}\label{eq:rescalerules}
d^{2}y\,d^{2}y'=\xi^{4}\,d^{2}w\,d^{2}w'\,,\qquad
K(y,y',z)=\frac{K(w,w',z)}{\xi^{2}}\,,\qquad
e^{-ik(y'-y)\frac1\xi}=e^{-ik(w'-w)}\ \ \text{\emph{exactly}, since } y'-y=\xi\,(w'-w)\,.
\end{align}
At fixed $(w,w',z)$ the square bracket of \eqref{eq:combo} behaves as
\begin{align}
\mathcal K(x-y)\!\cdot\!\mathcal K(x'-y')\,W(y,y')-K(x,x',z)\,W(z,z)
\;\xrightarrow[\ \xi\to0\ ]{}\;
\mathcal K(x-z)\!\cdot\!\mathcal K(x'-z)\,W(z,z)-K(x,x',z)\,W(z,z)=0\,,
\end{align}
because $y,y'\to z$ and $\mathcal K(x-z)\cdot\mathcal K(x'-z)=K(x,x',z)$ identically: the bracket
vanishes \emph{linearly} in $\xi$ (the kernels and the Wilson lines are smooth at $z$). Moreover,
in the variables $(w,w',z)$ the $z$-integral of \eqref{eq:combo} converges at large $z$ --- the
parent points $y,y'$ travel with $z$, so the two terms of the bracket fall off together --- which
is why the combination can be evaluated directly in $d=2$: it contains no $\frac{2}{\epsilon}$, no
$\gamma_{E}$ and no $\mu$. And because the bracket is $O(\xi)$, the soft region contributes
$\int_{\bar\xi}\frac{d\xi}{\xi}\times O(\xi)=\text{finite}$: no $\log\frac{\vee}{k^{+}}$ survives
in the combination. All the ultraviolet data and the entire rapidity logarithm of the soft channel
are carried by the subtraction term alone, which we now evaluate.

\textbf{Step 4: the exact value of the subtraction.} Eq.~\eqref{eq:sub} depends on $\xi$ only
through the phase, so the same substitution \eqref{eq:rescale} computes it in closed form: the
Jacobian $\xi^{4}$ times the kernel rescaling $\xi^{-2}$ gives
$\frac{d^{3}N^{\rm sub}_{\xi}}{d^{3}k}=\xi^{2}\times(\xi\text{-independent integral})$, and the
$\xi^{2}$ cancels the explicit $\frac{1}{\xi^{2}}$ of the DGLAP measure, leaving only
\begin{align}\label{eq:softint}
\int_{\bar\xi}^{1}d\xi\,\frac{1-\xi}{\xi}
=\Big[\log\xi-\xi\Big]_{\xi=\bar\xi}^{\,\xi=1}
=-\log\bar\xi-(1-\bar\xi)
=\log\frac{\vee}{k^{+}}-\Big(1-\frac{k^{+}}{\vee}\Big)
\end{align}
(upper minus lower again; the result is positive, as it must be for a positive integrand). Hence
\begin{align}\label{eq:subvalue}
-\frac{g^{2}}{8\pi^{2}}\int_{\bar\xi}^{1}\frac{d\xi}{\xi^{2}}\,2N_c\,\frac{1-\xi}{\xi}\,\frac{d^{3}N^{\rm sub}_{\xi}}{d^{3}k} =-\frac{1}{(2\pi)^{3}}\frac{g^{4}N_{c}}{8\pi^{5}}\frac{1}{k^{+}}\Big[\log\frac{\vee}{k^{+}}-\Big(1-\frac{k^{+}}{\vee}\Big)\Big]\int_{x,x',z,w,w'}e^{-ik\cdot(w'-w)}K(w,w',z)\,K(x,x',z)\,W(z,z)\,.
\end{align}
\textbf{Step 5: add and subtract.} The subtraction term was \emph{added} to the cross-section as
the ``soft subtraction'' entry of the finite block \eqref{eq:finite} (that entry is verbatim the
left-hand side of \eqref{eq:subvalue}; it is what renders the soft channel of the DGLAP term
safe via \eqref{eq:combo}), so its exact value \eqref{eq:subvalue} must be \emph{subtracted} back
to change nothing. Since \eqref{eq:subvalue} carries an overall minus sign, subtracting it back
means \emph{adding}
$+\frac{1}{(2\pi)^{3}}\frac{g^{4}N_{c}}{8\pi^{5}}\frac{1}{k^{+}}
\big[\log\frac{\vee}{k^{+}}-\big(1-\frac{k^{+}}{\vee}\big)\big]\int K\,K\,W(z,z)$: the
$\log\frac{\vee}{k^{+}}$ part (overall $+$) is the
\textbf{BFKL contribution of this row},
\begin{align}\label{eq:bfkl}
\boxed{\;\frac{d^{3}N^{\rm BFKL}_{\rm G\text{-}II,\,R\,IV}}{d^{3}k} = +\frac{1}{(2\pi)^{3}}\,\frac{g^{4}\,N_{c}}{8\pi^{5}}\,\frac{1}{k^{+}}\,{\color{blue}\log\Big(\frac{\vee}{k^{+}}\Big)}\int_{x,x',z,w,w'}e^{-ik\cdot (w'-w)}\,K(w,w',z)\,K(x,x',z)\,
\left[U^{ce'}(z) - U^{ce'}(x')\right]\left[U^{ce}(z) - U^{ce}(x)\right]\rho^{e'}(x')\rho^{e}(x)\;}
\end{align}
moved to the BFKL file (it appears there verbatim); its coefficient is manifestly UV finite and
$\mu$-independent, as a BFKL kernel must be. The $-\big(1-\frac{k^{+}}{\vee}\big)$ part remains as
the boundary term below.

\section{Result: the finite terms}
Collecting everything, Eq.~\eqref{eq:start} decomposes as
\begin{align}
\frac{d^{3}N_{\rm NLO}}{d^{3}k}
=\frac{d^{3}N^{\rm BFKL}}{d^{3}k}\Big|_{\eqref{eq:bfkl}}
+\frac{d^{3}N^{\rm JIMWLK}}{d^{3}k}\Big|_{\eqref{eq:jimwlk}}
+\frac{d^{3}N^{\rm finite}}{d^{3}k}\,,
\end{align}
with the finite part containing the DGLAP and RC terms:
\begin{align}\label{eq:finite}
    \frac{d^{3}N^{\rm finite}}{d^{3}k}  & =
   \underbrace{{\color{violet} -\,\frac{1}{(2\pi)^{3}}\,\frac{g^{2}}{2\pi^{2}}\,\frac{g^{2}}{8\pi^{2}}\,\frac{1}{k^{+}}\,\int_{k^{+}/\vee}^{1}\, \frac{d\xi}{\xi^{2}}\,\mathcal{P}_{gg}(\xi)\,\int_{x,y,x',y'}\, e^{-ik\cdot (y'-y)\,\frac{1}{\xi}}\, \KK\,\Leps\,\Wbr}}_{\textstyle \rm DGLAP}
    \notag \\ & \underbrace{-\,\frac{1}{(2\pi)^{3}}\,\frac{g^{4}\,N_{c}}{8\pi^{5}}\, \frac{1}{k^{+}}\,\int_{k^{+}/\vee}^{1}\,\frac{d\xi}{\xi^{2}}\,\frac{(1-\xi)}{\xi}\,\int_{x,y,x',y',z}\, e^{-ik\cdot (y'-y)\,\frac{1}{\xi}}\,K(y,y',z)\,K(x,x',z)\,
\left[U^{ce'}(z) - U^{ce'}(x')\right]\left[U^{ce}(z) - U^{ce}(x)\right]\rho^{e'}(x')\,\rho^{e}(x)}_{\textstyle \rm soft\ subtraction} \notag \\ & \underbrace{{\color{green!60!black}+\frac{1}{(2\pi)^{3}}\, \frac{g^{2}}{2\pi^{2}}\,\frac{g^{2}\,\beta_{0}}{16\pi^{2}}\, \frac{1}{k^{+}}\,\int_{x,y,x',y'}\, e^{-ik\cdot (y'-y)}\, \KK\;\Leps\;\Wbr}}_{\textstyle \rm RC}\notag \\ &
    \underbrace{-\frac{1}{(2\pi)^{3}}\,\frac{g^{4}\,N_{c}}{8\pi^{5}}\,\frac{1}{k^{+}}\,\Big(1-\frac{k^{+}}{\vee}\Big)\int_{x,x',z,w,w'}e^{-ik\cdot (w'-w)}\,K(w,w',z)\,K(x,x',z)\,
\left[U^{ce'}(z) - U^{ce'}(x')\right]\left[U^{ce}(z) - U^{ce}(x)\right]\rho^{e'}(x')\rho^{e}(x)}_{\textstyle \rm boundary} \notag \\
&  \underbrace{+\,\frac{1}{(2\pi)^{3}}\,\frac{g^{4}N_{c}}{8\pi^{4}}\, \frac{1}{k^{+}}\,\int_{\frac{k^{+}}{\vee}}^{1}\,d\xi\, \frac{(1-\xi)}{\xi}\, \int_{x,y,x',y'}\,e^{-ik\cdot (y'-y)\frac{1}{\xi}}\,\Bigg[\KK -2\,\frac{(x-y)\cdot (y'-y)\, (x'-y')\cdot (y'-y)}{(x-y)^{2}\,(x'-y')^{2}\,(y'-y)^{2}}\Bigg]\Wbr}_{\textstyle \rm azimuthal}\,,
\end{align}
with $\mathcal P_{gg}$ in the \emph{standard} $[0,1]$ plus-prescription convention
\eqref{eq:stdplus}--\eqref{eq:pgg}, $\Leps=\frac{2}{\epsilon}+2\gamma_{E}
+\log\frac{(y'-y)^{2}\mu^{2}_{\overline{MS}}}{4}$, and $d_{\perp}\to2$ set in the azimuthal term
(no pole multiplies it).

\subsection*{Bookkeeping}
\begin{itemize}
\item[(i)] (soft subtraction) $+$ (boundary) $+$ (BFKL entry \eqref{eq:bfkl}) $=0$ identically ---
the construction is an exact add-and-subtract of \eqref{eq:subvalue}; all three objects are UV
finite. The physical origin of the BFKL logarithm is the $\xi\to0$ endpoint of the DGLAP term.
\item[(ii)] The $\frac{\beta_{0}}{2}\delta(1-\xi)$ piece inside $\mathcal P_{gg}$ cancels the RC term
identically; the soft $2N_{c}\frac{1-\xi}{\xi}$ piece of the DGLAP term combines with the
subtraction into the $\mu$-free form \eqref{eq:combo}. The only surviving ultraviolet logarithms
are $\Leps\otimes 2N_{c}\big[\frac{\xi}{(1-\xi)_{+}}+\xi(1-\xi)\big]$ --- the physical collinear
(DGLAP) kernel of the measured gluon, including its $(1-\xi)_{+}$ endpoint
$\log(1-\frac{k^{+}}{\vee})$ piece --- and the RC term; these enter the final pole accounting
(coupling renormalization with $\beta_{0}=\frac{11N_{c}}{3}$ together with the Group-I ultraviolet
terms, plus the collinear subtraction for the produced gluon).
\item[(iii)] Together with \eqref{eq:jimwlk} and \eqref{eq:bfkl}, Eq.~\eqref{eq:finite} reproduces
the starting expression \eqref{eq:start} identically, up to $O(\Lambda)$ (endpoint recombination
and extension of regular $\xi$-integrals to $\xi=1$) and $O(\epsilon)\times{\rm finite}$.
\end{itemize}
\end{document}
